\chapter{Thiết kế phần mềm}

\section{Tổng quan thiết kế phần mềm}

Thiết kế phần mềm của hệ thống được chia thành ba phần: thiết kế dữ liệu, thiết kế xử lý và thiết kế giao diện. Thiết kế dữ liệu mô tả cách hệ thống lưu trữ và ràng buộc dữ liệu; thiết kế xử lý mô tả luồng nghiệp vụ, API và business rules; thiết kế giao diện mô tả cách người dùng tương tác với hệ thống.

\section{Thiết kế dữ liệu}

\subsection{Mục tiêu thiết kế dữ liệu}

Cơ sở dữ liệu được thiết kế nhằm đảm bảo dữ liệu bệnh nhân, lượt khám, phiếu khám, đơn thuốc, hóa đơn, thanh toán, lịch hẹn, xét nghiệm và kho thuốc được lưu trữ nhất quán. Hệ thống sử dụng PostgreSQL và Prisma ORM. Schema hiện có 37 models, 12 enums và 3 migrations.

\subsection{ERD tổng thể và cách tách sơ đồ}

Hệ thống có 37 models chia thành Phase 1 (20 models lõi) và Phase 2 (17 models mở rộng). ERD Phase 1 và ERD Phase 2 Delta được trình bày trên trang ngang để đảm bảo chữ đủ lớn và dễ đọc.

\LandscapeFigure{figures/ch03-software-design/erd-phase1.pdf}{ERD Phase 1 -- nghiệp vụ lõi phòng mạch}{fig:erd-phase1}
\LandscapeFigure{figures/ch03-software-design/erd-phase2-delta.pdf}{ERD Phase 2 Delta -- các model mở rộng}{fig:erd-phase2-delta}

\subsection{Các nhóm model chính}

\begin{longtable}{|p{3.3cm}|p{5.5cm}|p{5.4cm}|}
\caption{Các nhóm model trong cơ sở dữ liệu}\\
\hline
\textbf{Nhóm} & \textbf{Models} & \textbf{Vai trò thiết kế} \\
\hline
Identity \& Access & User, Role, Permission, UserRole, RolePermission, RefreshToken, AuditLog & Quản lý tài khoản, vai trò, quyền, token và nhật ký hệ thống \\
\hline
Patients \& Visits & Patient, Visit & Quản lý hồ sơ bệnh nhân, lượt khám và số thứ tự \\
\hline
Examination \& Prescription & Examination, Disease, Diagnosis, Drug, Prescription, PrescriptionItem & Lưu phiếu khám, chẩn đoán, danh mục bệnh, thuốc và đơn thuốc \\
\hline
Billing \& Regulation & Invoice, InvoiceItem, Payment, RegulationVersion, RegulationItem & Quản lý hóa đơn, thanh toán và quy định phòng mạch \\
\hline
Organization & Department, Room, DoctorProfile, StaffSchedule & Quản lý tổ chức phòng khám, phòng và bác sĩ \\
\hline
Appointment \& Queue & Appointment, QueueTicket & Quản lý lịch hẹn và hàng đợi \\
\hline
Vitals \& Services \& Lab & VitalSign, ServiceCatalog, LabTestCatalog, ServiceOrder, LabOrder, LabSample, LabResult & Sinh hiệu, chỉ định dịch vụ và quy trình xét nghiệm \\
\hline
Inventory \& Pharmacy & StockLot, StockMovement, Dispense, DispenseItem & Quản lý tồn kho, biến động kho và cấp phát thuốc \\
\hline
\end{longtable}

\subsection{Enum và state machine}

\begin{longtable}{|p{4cm}|p{8cm}|p{3cm}|}
\caption{Danh sách 12 enum trong database}\\
\hline
\textbf{Enum} & \textbf{Values} & \textbf{Model sử dụng} \\
\hline
UserStatus & ACTIVE, INACTIVE, LOCKED & User \\
\hline
VisitStatus & REGISTERED, WAITING, IN\_EXAMINATION, COMPLETED, CANCELLED & Visit \\
\hline
ExaminationStatus & OPEN, COMPLETED, CANCELLED & Examination \\
\hline
InvoiceStatus & DRAFT, ISSUED, PARTIALLY\_PAID, PAID, VOID & Invoice \\
\hline
PaymentMethod & CASH, TRANSFER, CARD & Payment \\
\hline
AppointmentStatus & SCHEDULED, CHECKED\_IN, CANCELLED, NO\_SHOW & Appointment \\
\hline
QueueStatus & WAITING, CALLED, IN\_SERVICE, DONE, SKIPPED, CANCELLED & QueueTicket \\
\hline
ServiceType & CONSULTATION, LAB\_TEST, PROCEDURE & ServiceCatalog \\
\hline
DispenseStatus & PENDING, DISPENSED, CANCELLED & Dispense \\
\hline
StockMovementType & IN, OUT, ADJUSTMENT & StockMovement \\
\hline
LabOrderStatus & ORDERED, SAMPLE\_COLLECTED, RESULT\_ENTERED, VERIFIED, CANCELLED & LabOrder \\
\hline
ServiceOrderStatus & ORDERED, IN\_PROGRESS, COMPLETED, CANCELLED & ServiceOrder \\
\hline
\end{longtable}

\begin{longtable}{|p{2.6cm}|p{3.5cm}|p{4.8cm}|p{3.4cm}|}
\caption{Chuyển trạng thái các entity lõi trong hệ thống}\label{tab:state-transitions}\\
\hline
\textbf{Entity} & \textbf{Trạng thái hiện tại} & \textbf{Sự kiện kích hoạt} & \textbf{Trạng thái mới} \\
\hline
\multicolumn{4}{|l|}{\textbf{Visit}} \\
\hline
Visit & -- & Lễ tân tạo lượt khám & WAITING \\
\hline
Visit & WAITING & Bác sĩ mở phiên khám & IN\_EXAMINATION \\
\hline
Visit & IN\_EXAMINATION & Phiếu khám hoàn tất & COMPLETED \\
\hline
Visit & WAITING / IN\_EXAMINATION & Hủy lượt khám & CANCELLED \\
\hline
\multicolumn{4}{|l|}{\textbf{Examination}} \\
\hline
Examination & -- & Bác sĩ mở phiên khám & OPEN \\
\hline
Examination & OPEN & Hoàn tất phiếu khám & COMPLETED \\
\hline
Examination & OPEN & Hủy & CANCELLED \\
\hline
\multicolumn{4}{|l|}{\textbf{Invoice}} \\
\hline
Invoice & -- & Thu ngân lập hóa đơn & ISSUED \\
\hline
Invoice & ISSUED & Thanh toán chưa đủ & PARTIALLY\_PAID \\
\hline
Invoice & ISSUED / PARTIALLY\_PAID & Thanh toán đủ số tiền & PAID \\
\hline
Invoice & ISSUED / PARTIALLY\_PAID & Hủy hóa đơn & VOID \\
\hline
\multicolumn{4}{|l|}{\textbf{QueueTicket (Phase 2)}} \\
\hline
QueueTicket & -- & Check-in tạo vé hàng đợi & WAITING \\
\hline
QueueTicket & WAITING & Gọi tên bệnh nhân & CALLED \\
\hline
QueueTicket & CALLED & Bắt đầu phục vụ & IN\_SERVICE \\
\hline
QueueTicket & IN\_SERVICE & Hoàn tất & DONE \\
\hline
QueueTicket & WAITING / CALLED & Hủy hoặc bỏ qua & CANCELLED \\
\hline
\end{longtable}

Bảng~\ref{tab:state-transitions} tóm tắt các chuyển trạng thái được phép trong quy trình vận hành. Ví dụ, \code{VisitStatus} đi từ WAITING sang IN\_EXAMINATION khi bác sĩ mở phiên khám, sau đó sang COMPLETED khi phiếu khám hoàn tất; các trạng thái CANCELLED/VOID chỉ được dùng cho luồng hủy hoặc vô hiệu hóa có kiểm soát. Việc định nghĩa rõ chuyển trạng thái giúp backend tránh cập nhật tùy tiện và giúp frontend hiển thị hành động phù hợp theo từng trạng thái.

\subsection{Ràng buộc dữ liệu quan trọng}

\begin{longtable}{|p{4.2cm}|p{3.2cm}|p{7cm}|}
\caption{Ràng buộc dữ liệu quan trọng}\\
\hline
\textbf{Ràng buộc} & \textbf{Model} & \textbf{Ý nghĩa} \\
\hline
UUID primary key & Tất cả model & Không dùng serial integer, giảm khả năng đoán ID \\
\hline
\code{patientCode @unique} & Patient & Mỗi bệnh nhân có mã riêng \\
\hline
\code{citizenId @unique} & Patient & Tránh trùng định danh bệnh nhân \\
\hline
\code{@@unique([patientId, visitDate])} & Visit & Ngăn một bệnh nhân có hai lượt khám trong cùng ngày \\
\hline
\code{@@unique([visitDate, queueNumber])} & Visit & Số thứ tự không trùng trong ngày \\
\hline
\code{visitId @unique} & Examination, Invoice & Một lượt khám có một phiếu khám và một hóa đơn chính \\
\hline
\code{examinationId @unique} & Prescription & Một phiếu khám có tối đa một đơn thuốc \\
\hline
\code{@@unique([prescriptionId, drugId])} & PrescriptionItem & Không kê trùng cùng một thuốc trong một đơn \\
\hline
\code{@@unique([drugId, lotNumber])} & StockLot & Không trùng mã lô cho cùng một thuốc \\
\hline
\end{longtable}

\subsection{Quyết định thiết kế dữ liệu}

\begin{table}[H]
\centering
\caption{Các quyết định thiết kế dữ liệu quan trọng}
\begin{tabularx}{\textwidth}{|p{4cm}|X|}
\hline
\textbf{Quyết định} & \textbf{Lý do} \\
\hline
Dùng UUID làm primary key & Phù hợp khi hệ thống mở rộng, không để lộ thứ tự tăng dần của bản ghi \\
\hline
Soft-delete bằng \code{isActive} & Giữ lịch sử đối với Disease, Drug, ServiceCatalog thay vì xóa cứng \\
\hline
Snapshot giá trong PrescriptionItem & Lưu đơn giá tại thời điểm kê đơn, tránh thay đổi lịch sử khi giá thuốc thay đổi \\
\hline
RegulationVersion + RegulationItem & Cho phép quản lý phiên bản quy định và chỉ kích hoạt một phiên bản tại một thời điểm \\
\hline
StockLot + StockMovement & Phục vụ quản lý tồn kho, hạn dùng và audit biến động kho \\
\hline
\end{tabularx}
\end{table}

\section{Thiết kế xử lý}

\subsection{Nguyên tắc thiết kế xử lý}

Các xử lý nghiệp vụ quan trọng được đặt tại service layer của backend. Frontend có thể kiểm tra nhập liệu để cải thiện trải nghiệm người dùng, nhưng không thay thế kiểm tra nghiệp vụ. Nhờ vậy, dù API được gọi từ frontend hay từ Swagger/Postman, các quy tắc như chống trùng lượt khám, không thanh toán vượt số tiền còn lại, không cấp phát thuốc khi không đủ tồn vẫn được kiểm soát ở backend.

\subsection{API inventory tổng quan}

Hệ thống có 92 endpoints, trong đó 41 endpoints thuộc Phase 1 và 51 endpoints thuộc Phase 2. API prefix là \code{/api/v1}; Swagger UI đặt tại \code{http://localhost:3000/api/docs}.

\begin{table}[H]
\centering
\caption{Tổng quan endpoint theo nhóm module}
\begin{tabularx}{\textwidth}{|p{4.2cm}|p{2.5cm}|X|}
\hline
\textbf{Nhóm} & \textbf{Số endpoint} & \textbf{Vai trò} \\
\hline
Auth, Users, RBAC & 13 & Đăng nhập, token, tài khoản, phân quyền \\
\hline
Patients, Visits, Examinations & 13 & Hồ sơ bệnh nhân, lượt khám, phiếu khám, kê đơn \\
\hline
Billing, Diseases, Drugs, Regulations, Reports & 15 & Hóa đơn, thanh toán, danh mục, quy định, báo cáo \\
\hline
Appointments, Queue, Vitals & 12 & Lịch hẹn, hàng đợi, sinh hiệu \\
\hline
Services, Lab & 14 & Dịch vụ và xét nghiệm \\
\hline
Inventory, Pharmacy, Organization, Audit & 25 & Kho thuốc, cấp phát, tổ chức, audit log \\
\hline
\end{tabularx}
\end{table}

\subsection{Ma trận phân quyền RBAC}

Backend sử dụng \code{JwtAuthGuard}, \code{RolesGuard} và decorator \code{@Roles} để kiểm soát quyền theo nhóm endpoint. Bảng dưới đây tóm tắt các quyền chính theo vai trò Phase 1 và một số endpoint mở rộng.

\begin{longtable}{|p{5cm}|c|c|c|c|c|}
\caption{Ma trận phân quyền theo nhóm endpoint chính}\\
\hline
\textbf{Nhóm endpoint} & \textbf{ADMIN} & \textbf{RECEPT.} & \textbf{DOCTOR} & \textbf{CASHIER} & \textbf{MANAGER} \\
\hline
POST /auth/login & Yes & Yes & Yes & Yes & Yes \\
\hline
GET/POST /users & Yes & -- & -- & -- & -- \\
\hline
GET/POST /patients & Yes & Yes & Yes & -- & -- \\
\hline
POST /visits & Yes & Yes & -- & -- & -- \\
\hline
GET /visits & Yes & Yes & Yes & -- & Yes \\
\hline
POST /visits/:id/open-examination & Yes & -- & Yes & -- & -- \\
\hline
PATCH /examinations/:id & Yes & -- & Yes & -- & -- \\
\hline
POST /invoices & Yes & -- & -- & Yes & -- \\
\hline
GET/POST /invoices/:id/payments & Yes & -- & -- & Yes & -- \\
\hline
GET /reports/monthly & Yes & -- & -- & -- & Yes \\
\hline
GET/POST /regulations & Yes & -- & -- & -- & Yes \\
\hline
GET /audit & Yes & -- & -- & -- & Yes \\
\hline
\end{longtable}

\subsection{Danh sách endpoint Phase 1 tiêu biểu}

Bảng sau liệt kê các endpoint Phase 1 quan trọng nhất để chứng minh thiết kế API có thể truy vết trực tiếp từ yêu cầu và use case.

\begin{longtable}{|p{2cm}|p{5.2cm}|p{3cm}|p{4.2cm}|}
\caption{Endpoint Phase 1 tiêu biểu}\\
\hline
\textbf{Method} & \textbf{Path} & \textbf{Role chính} & \textbf{Use case liên quan} \\
\hline
POST & /api/v1/auth/login & Tất cả & UC01 đăng nhập \\
\hline
POST & /api/v1/auth/refresh & Tất cả & Duy trì phiên đăng nhập \\
\hline
POST & /api/v1/auth/logout & Tất cả & Thu hồi refresh token \\
\hline
GET & /api/v1/auth/me & Tất cả & Bootstrap phiên người dùng \\
\hline
GET & /api/v1/users & ADMIN & UC02 quản lý tài khoản \\
\hline
POST & /api/v1/users & ADMIN & UC02 quản lý tài khoản \\
\hline
GET & /api/v1/patients & RECEPTIONIST, DOCTOR & UC04 tra cứu hồ sơ \\
\hline
POST & /api/v1/patients & RECEPTIONIST & UC05 tạo hồ sơ \\
\hline
GET & /api/v1/patients/:id/medical-history & DOCTOR & UC12 xem lịch sử khám \\
\hline
POST & /api/v1/visits & RECEPTIONIST & UC06/UC07 tạo lượt khám \\
\hline
GET & /api/v1/visits & RECEPTIONIST, DOCTOR & UC07 xem danh sách lượt khám \\
\hline
POST & /api/v1/visits/:id/open-examination & DOCTOR & UC08 mở phiên khám \\
\hline
PATCH & /api/v1/examinations/:id & DOCTOR & UC09 ghi phiếu khám \\
\hline
POST & /api/v1/examinations/:id/complete & DOCTOR & UC13 hoàn tất khám \\
\hline
POST & /api/v1/invoices & CASHIER & UC14 lập hóa đơn \\
\hline
POST & /api/v1/invoices/:id/payments & CASHIER & UC15 ghi nhận thanh toán \\
\hline
GET & /api/v1/reports/monthly & MANAGER & UC18 báo cáo tháng \\
\hline
\end{longtable}

\subsection{DTO và response shape tiêu biểu}

Các request quan trọng được mô hình hóa bằng DTO để kiểm soát dữ liệu đầu vào. Ví dụ, DTO tạo lượt khám tối thiểu cần \code{patientId} và ngày khám; response trả về \code{id}, \code{patientId}, \code{visitDate}, \code{queueNumber} và \code{status}. DTO thanh toán cần \code{amount}, \code{method} và ghi chú tùy chọn; response trả về hóa đơn đã cập nhật kèm danh sách payment. Cách thiết kế này giúp frontend không phụ thuộc trực tiếp vào schema database và giúp backend kiểm tra dữ liệu ở biên API.

\subsection{Business rules trọng tâm}

\begin{longtable}{|p{1.8cm}|p{3cm}|p{7cm}|p{2.8cm}|}
\caption{Business rules trọng tâm}\\
\hline
\textbf{BR} & \textbf{Module} & \textbf{Mô tả} & \textbf{Status} \\
\hline
BR-01 & auth & Validate email/password, sai thông tin trả 401 & CONFIRMED \\
\hline
BR-04 & patients & Không tạo bệnh nhân có citizenId trùng & CONFIRMED \\
\hline
BR-05 & visits & Không tạo hai lượt khám cùng bệnh nhân trong cùng ngày & CONFIRMED \\
\hline
BR-06 & visits & Kiểm tra quota lượt khám trong ngày theo regulation & CONFIRMED \\
\hline
BR-07 & visits & Sinh queue number trong transaction & CONFIRMED \\
\hline
BR-10 & examinations & Chỉ phiếu khám OPEN mới được cập nhật & CONFIRMED \\
\hline
BR-12 & examinations & Không hoàn tất khám nếu chưa có diagnosis & CONFIRMED \\
\hline
BR-14 & examinations & Chưa chặn complete khi còn pending lab/service order & RISK \\
\hline
BR-15 & billing & Chỉ lập invoice khi Visit đã COMPLETED & CONFIRMED \\
\hline
BR-17 & billing & Không thanh toán vượt số tiền còn lại & CONFIRMED \\
\hline
BR-20 & regulations & Activate regulation version trong transaction & CONFIRMED \\
\hline
BR-21--24 & pharmacy & Kiểm tra hết hạn, tồn kho, duplicate dispense, transaction trừ kho & CONFIRMED \\
\hline
BR-25 & lab & Lab flow ORDERED, SAMPLE\_COLLECTED, RESULT\_ENTERED, VERIFIED & CONFIRMED \\
\hline
BR-26 & vitals & Tự động tính BMI từ cân nặng và chiều cao & CONFIRMED \\
\hline
BR-28 & queue & State machine hàng đợi WAITING, CALLED, IN\_SERVICE, DONE/SKIPPED & CONFIRMED \\
\hline
\end{longtable}

\subsection{Các luồng xử lý chính}

\begin{longtable}{|p{0.5cm}|p{3.2cm}|p{5.5cm}|p{4.6cm}|}
\caption{Các bước trong luồng tạo lượt khám (BR-05, BR-06, BR-07)}\label{tab:visit-sequence}\\
\hline
\textbf{\#} & \textbf{Actor/Component} & \textbf{Hành động} & \textbf{Kết quả / Điều kiện} \\
\hline
1 & Receptionist $\rightarrow$ Frontend & Chọn bệnh nhân, nhập ngày khám, submit & POST /visits \{patientId, visitDate\} \\
\hline
2 & VisitsService $\rightarrow$ DB & Tìm bệnh nhân theo patientId & 404 Not Found nếu không tồn tại \\
\hline
3 & VisitsService $\rightarrow$ DB & Kiểm tra lượt khám trùng cùng bệnh nhân và ngày & 409 Conflict nếu đã có lượt khám hôm nay \\
\hline
4 & VisitsService $\rightarrow$ RegulationsService & Lấy quota lượt khám tối đa trong ngày & Giá trị từ regulation hiện hành \\
\hline
5 & VisitsService $\rightarrow$ DB & Đếm lượt khám trong ngày & 409 Conflict nếu đã đạt quota \\
\hline
6 & VisitsService $\rightarrow$ DB & SELECT MAX(queueNumber) + 1 cho ngày khám & Số thứ tự tiếp theo \\
\hline
7 & VisitsService $\rightarrow$ DB & INSERT Visit trong \code{\$transaction} (BR-05, BR-07) & Visit mới trạng thái WAITING \\
\hline
8 & VisitsController $\rightarrow$ Frontend & 201 Created \{id, queueNumber, status\} & Hiển thị số thứ tự cho lễ tân \\
\hline
\end{longtable}

Luồng tạo lượt khám bắt đầu khi lễ tân chọn bệnh nhân và gửi yêu cầu. Backend lần lượt kiểm tra bệnh nhân tồn tại, chưa có lượt khám trùng ngày, chưa vượt quota và tạo số thứ tự trong transaction. Nếu vi phạm bất kỳ điều kiện nào, backend trả lỗi nghiệp vụ 409/404 cho frontend hiển thị.

\begin{longtable}{|p{0.5cm}|p{3cm}|p{5cm}|p{5.2cm}|}
\caption{Các bước trong luồng khám bệnh và kê đơn (BR-10, BR-12)}\label{tab:examination-sequence}\\
\hline
\textbf{\#} & \textbf{Actor/Component} & \textbf{Hành động} & \textbf{Kết quả / Điều kiện} \\
\hline
\multicolumn{4}{|l|}{\textbf{Giai đoạn 1 -- Mở phiên khám}} \\
\hline
1 & Doctor $\rightarrow$ Frontend & Chọn lượt khám WAITING, mở phiên & POST /visits/\{id\}/open-examination \\
\hline
2 & ExaminationsService $\rightarrow$ DB & Kiểm tra trạng thái Visit hợp lệ & 400 Bad Request nếu sai trạng thái \\
\hline
3 & ExaminationsService $\rightarrow$ DB & INSERT Examination OPEN; UPDATE Visit IN\_EXAMINATION & 201 Created, form khám mở ra \\
\hline
\multicolumn{4}{|l|}{\textbf{Giai đoạn 2 -- Ghi chẩn đoán}} \\
\hline
4 & Doctor $\rightarrow$ Frontend & Nhập triệu chứng, chẩn đoán, bệnh từ danh mục, phí & PATCH /examinations/\{id\} \\
\hline
5 & ExaminationsService $\rightarrow$ DB & Kiểm tra Examination chưa COMPLETED (BR-10) & 400 Bad Request nếu đã hoàn tất \\
\hline
6 & ExaminationsService $\rightarrow$ DB & UPDATE Examination & 200 OK \\
\hline
\multicolumn{4}{|l|}{\textbf{Giai đoạn 3 -- Kê đơn thuốc}} \\
\hline
7 & Doctor $\rightarrow$ Frontend & Thêm thuốc và số lượng vào đơn & PATCH /examinations/\{id\} (prescription items) \\
\hline
8 & PrescriptionsService $\rightarrow$ DB & Validate thuốc active, không trùng trong đơn & 400 Bad Request nếu vi phạm \\
\hline
9 & PrescriptionsService $\rightarrow$ DB & Transaction upsert Prescription + Items với snapshot giá & 200 OK, đơn thuốc lưu xong \\
\hline
\multicolumn{4}{|l|}{\textbf{Giai đoạn 4 -- Hoàn tất khám}} \\
\hline
10 & Doctor $\rightarrow$ Frontend & Nhấn Hoàn tất & POST /examinations/\{id\}/complete \\
\hline
11 & ExaminationsService $\rightarrow$ DB & Kiểm tra có chẩn đoán tối thiểu (BR-12) & 400 Bad Request nếu thiếu chẩn đoán \\
\hline
12 & ExaminationsService $\rightarrow$ DB & UPDATE Exam COMPLETED; UPDATE Visit COMPLETED & 200 OK, chuyển sang thu ngân \\
\hline
\end{longtable}

Luồng khám bệnh gồm 4 giai đoạn tuần tự: mở phiên khám, ghi chẩn đoán, kê đơn và hoàn tất. Hệ thống kiểm tra trạng thái phiếu khám, thuốc active và chẩn đoán tối thiểu trước mỗi bước quan trọng.

\begin{longtable}{|p{0.5cm}|p{3.2cm}|p{5.5cm}|p{4.6cm}|}
\caption{Các bước trong luồng lập hóa đơn và thanh toán (BR-15, BR-17)}\label{tab:payment-sequence}\\
\hline
\textbf{\#} & \textbf{Actor/Component} & \textbf{Hành động} & \textbf{Kết quả / Điều kiện} \\
\hline
\multicolumn{4}{|l|}{\textbf{Giai đoạn 1 -- Lập hóa đơn}} \\
\hline
1 & Cashier $\rightarrow$ Frontend & Mở lượt khám đã COMPLETED & POST /invoices \{visitId\} \\
\hline
2 & BillingService $\rightarrow$ DB & Tải Visit + Examination + PrescriptionItems & Dữ liệu phí khám và thuốc \\
\hline
3 & BillingService & Kiểm tra Visit trạng thái COMPLETED (BR-15) & 400 Bad Request nếu chưa hoàn tất \\
\hline
4 & BillingService $\rightarrow$ DB & Tính InvoiceItems, INSERT Invoice + Items & 201 Created, hóa đơn ISSUED \\
\hline
\multicolumn{4}{|l|}{\textbf{Giai đoạn 2 -- Ghi nhận thanh toán}} \\
\hline
5 & Cashier $\rightarrow$ Frontend & Nhập số tiền và phương thức thanh toán & POST /invoices/\{id\}/payments \\
\hline
6 & BillingService $\rightarrow$ DB & Tải Invoice, tính remaining = total -- paid & Số tiền còn lại \\
\hline
7 & BillingService & Kiểm tra amount > 0 & 400 Bad Request nếu không dương \\
\hline
8 & BillingService & Kiểm tra amount $\leq$ remaining (BR-17) & 400 Bad Request nếu vượt số tiền còn lại \\
\hline
9 & BillingService $\rightarrow$ DB & Transaction: INSERT Payment, UPDATE paidAmount, UPDATE InvoiceStatus & 200 OK, trạng thái PARTIALLY\_PAID hoặc PAID \\
\hline
\end{longtable}

Thu ngân chỉ lập hóa đơn khi lượt khám đã COMPLETED. Khi ghi nhận thanh toán, hệ thống kiểm tra số tiền lớn hơn 0 và không vượt phần còn lại; mọi cập nhật thực hiện trong một transaction để đảm bảo nhất quán.

\begin{longtable}{|p{0.5cm}|p{3.2cm}|p{5.5cm}|p{4.6cm}|}
\caption{Các bước trong luồng cấp phát thuốc theo FEFO (BR-21--BR-24)}\label{tab:fefo-sequence}\\
\hline
\textbf{\#} & \textbf{Actor/Component} & \textbf{Hành động} & \textbf{Kết quả / Điều kiện} \\
\hline
1 & Pharmacist $\rightarrow$ Frontend & Mở đơn thuốc pending, chọn Phát thuốc & POST /pharmacy/dispense \{prescriptionId\} \\
\hline
2 & PharmacyService $\rightarrow$ DB & Tải Prescription + PrescriptionItems & Danh sách thuốc cần cấp phát \\
\hline
3 & InventoryService $\rightarrow$ DB & Với mỗi thuốc: truy vấn StockLot còn tồn, ORDER BY expiryDate ASC (FEFO) & Danh sách lô theo thứ tự hết hạn sớm nhất \\
\hline
4 & InventoryService & Kiểm tra tổng tồn kho $\geq$ số lượng cần & 400 Bad Request nếu không đủ tồn kho \\
\hline
5 & InventoryService & Phân bổ số lượng từ lô gần hết hạn nhất trước & Lot allocations theo FEFO \\
\hline
6 & PharmacyService $\rightarrow$ DB & Transaction: INSERT Dispense, DispenseItems; UPDATE StockLot.quantityOnHand; INSERT StockMovement OUT & 201 Created, tồn kho giảm tương ứng \\
\hline
\end{longtable}

\section{Thiết kế giao diện}

\subsection{Nguyên tắc thiết kế giao diện}

Giao diện được thiết kế theo hướng nghiệp vụ, trong đó mỗi vai trò chỉ nhìn thấy các chức năng cần thiết cho công việc của mình. Layout chính gồm sidebar điều hướng, topbar hiển thị thông tin người dùng và vùng nội dung chính. Các màn hình dữ liệu có trạng thái loading, empty, error và success.

\subsection{Danh sách màn hình chính}

\begin{longtable}{|p{4cm}|p{3.2cm}|p{7cm}|}
\caption{Danh sách màn hình giao diện chính}\\
\hline
\textbf{Màn hình} & \textbf{Actor} & \textbf{Mục đích} \\
\hline
LoginPage & Tất cả & Đăng nhập hệ thống \\
\hline
DashboardPage & Tất cả & Hiển thị tổng quan theo vai trò \\
\hline
PatientList/Create/Detail & RECEPTIONIST, DOCTOR, ADMIN & Quản lý hồ sơ bệnh nhân \\
\hline
VisitList/Create & RECEPTIONIST, DOCTOR & Quản lý lượt khám \\
\hline
ExaminationPage & DOCTOR & Lập phiếu khám, chẩn đoán, kê đơn, sinh hiệu, dịch vụ \\
\hline
InvoiceList/Detail & CASHIER, MANAGER & Quản lý hóa đơn, thanh toán \\
\hline
MonthlyReportPage & MANAGER & Báo cáo tháng, revenue breakdown \\
\hline
Catalog pages & ADMIN, MANAGER & Danh mục bệnh, thuốc, dịch vụ, quy định \\
\hline
Phase 2 pages & Nhiều vai trò & Lịch hẹn, queue, lab, inventory, pharmacy, organization, audit \\
\hline
\end{longtable}

\subsection{RBAC trên giao diện}

Frontend sử dụng \code{ProtectedRoute} để yêu cầu đăng nhập và \code{RequireRole} để kiểm tra role trước khi cho phép truy cập route. Sidebar cũng được lọc theo role để giảm thao tác sai. Tuy nhiên, frontend chỉ là lớp hỗ trợ trải nghiệm; backend vẫn là lớp kiểm soát quyền truy cập chính.

\section{Demo các chức năng trọng tâm}

\begin{longtable}{|p{3.2cm}|p{5.5cm}|p{5.5cm}|}
\caption{Demo thiết kế phần mềm}\\
\hline
\textbf{Nhóm demo} & \textbf{Nội dung trình bày} & \textbf{Minh chứng cần chèn} \\
\hline
Demo thiết kế dữ liệu & Quan hệ Patient, Visit, Examination, Prescription, Invoice & ERD Phase 1, bảng quan hệ chính \\
\hline
Demo ràng buộc dữ liệu & Chặn visit trùng cùng bệnh nhân trong ngày & Constraint + screenshot lỗi 409 \\
\hline
Demo thiết kế xử lý & Luồng tạo lượt khám qua frontend, controller, service, database & Sequence diagram tạo lượt khám \\
\hline
Demo xử lý thanh toán & Không cho thanh toán vượt số tiền còn lại & Sequence payment + API response lỗi \\
\hline
Demo giao diện & Login, dashboard, patient list, examination page, invoice page & 5 screenshot màn hình chính \\
\hline
\end{longtable}

\TwoDemoFigures
{figures/screenshots/patient-list.png}
{Danh sách bệnh nhân}
{figures/screenshots/invoice-detail.png}
{Chi tiết hóa đơn}
{Minh chứng giao diện cho hai nghiệp vụ quản lý bệnh nhân và thanh toán}
{fig:ui-demo-patient-invoice}
